% ====================================================================
% ABSTRACT (BAHASA INGGRIS)
% ====================================================================
% File ini berisi abstrak dalam bahasa Inggris (200-300 kata).
% Pastikan terjemahan akurat dan natural dalam bahasa Inggris.
% EDIT konten sesuai dengan penelitian Anda.
% ====================================================================

\frontmatterchapter{ABSTRACT}

% --- ISI ABSTRACT ---
% Ganti konten berikut dengan abstrak penelitian Anda dalam bahasa Inggris
\noindent The background of this research is \ldots (Explain the context and motivation of your research here, at least 2-3 sentences).

The research problem investigated in this study is \ldots (Explain the main problem to be solved).

The purpose of this research is to \ldots (Explain the objectives to be achieved from your research).

The research method used is \ldots (Explain the research method briefly: research type, sample, instruments, analysis techniques).

The research results show that \ldots (Explain the main findings from your research).

The conclusion is that \ldots (Provide a brief conclusion and implications of the research).

% --- ISI ALTERNATIF MENGGUNAKAN LIPSUM ---
% Jika Anda ingin placeholder lebih panjang, gunakan:
% \lipsum[8]

\vspace{1cm}

% --- KEYWORDS ---
\noindent \textbf{Keywords:} \textit{\keywordOne, \keywordTwo, \keywordThree}

\newpage

% ====================================================================
% PETUNJUK PENGISIAN ABSTRACT:
%
% 1. PANJANG: 200-300 kata (sesuai juknis)
% 2. STRUKTUR: Sama dengan abstrak Indonesia
%    - Background/Context
%    - Problem statement
%    - Research objectives
%    - Methodology (briefly)
%    - Main results
%    - Conclusions
%
% 3. GAYA PENULISAN:
%    - Gunakan formal academic English
%    - Past tense untuk metode
%    - Present tense untuk hasil
%    - Hindari kontraksi (don't → do not)
%
% 4. KEYWORDS:
%    - Sama dengan versi Indonesia (atau terjemahan yang sesuai)
%    - Edit di metadata.tex jika berbeda dari versi Indonesia
%
% 5. TIPS UMUM:
%    - Gunakan Grammarly atau editor online untuk validasi
%    - Hindari terjemahan literal dari bahasa Indonesia
%    - Baca ulang untuk natural flow dan grammar
%    - Konsultasikan dengan pembimbing untuk versi final
%
% 6. CATATAN PENTING:
%    - Abstract harus bisa berdiri sendiri (self-contained)
%    - Jangan mereferensikan ke bagian lain dokumen
%    - Hindari tabel, gambar, dan referensi
% ====================================================================
